%% THERMODYNAMIQUE %%
%
\section{Humidit\'e relative \`a partir de la temp\'erature du point
de ros\'ee sans seuillage}

\itemstyle{ Op\'erateur~: } {\sf HumW\_T\_Td\_no\_limit}
\index{HumW\_T\_Td\_no\_limit}

%%%%%OBJET%%%%%
\itemstyle{ Objet~: }
Cet op\'erateur calcule l'humidit\'e \`a partir de la temp\'erature 
du point de ros\'ee et de la temp\'erature de l'air, en utilisant les
formules relatives \`a l'eau liquide. 

%%%%%ARGUMENTS%%%%%
\bigskip
%%%%%ENTREES%%%%
\begin{tabarg}
\multicolumn{3}{l}{\titretab{ Entr\'ees}}\\
\hline
$T_s$ & Vecteur, {\tt float} & Temp\'erature statique de l'air ($K$ ou $^\circ C$) \\
$T_d$ & Vecteur, {\tt float} & Temp\'erature du point de ros\'ee ($K$ ou $^\circ C$) 
\end{tabarg}

\medskip
%%%%%SORTIES%%%%
\begin{tabarg}
\multicolumn{3}{l}{\titretab{Sorties}}\\
\hline
$H_u$ & Vecteur, {\tt float} & Humidit\'e relative ($\%$)\\
\end{tabarg}


%%%%%FORMULE%%%%
\itemstyle{ Formule }

%%%% THERMODYNAMIQUE
\[
H_u = 100 \times \frac{ e_w(T_d)}{e_w(T_s)} 
\]



%%%%COMMENTAIRE%%%%
\itemstyle{Commentaire~:}

La fonction $ew()$ est d\'ecrite en annexe
\ref{annexe-thermo}.

%%%%%AUTEUR%%%%%
\itemstyle{ Auteur~:}
CNRM/CAM/ITD

%%%%%REFERENCES%%%%%
\itemstyle{ R\'ef\'erences~:}
Triplet-Roche \cite{Triplet-Roche}
